\documentclass{IEEEtran}

% ======================
% Packages
% ======================
\usepackage{graphicx}
\usepackage{amsmath}
\usepackage{float}

\title{ECG Heartbeat Categorization}

\author{
    Pham Cong Duyet--23BI14136\\
    ML in Medicine -- Practical Work 1\\
    University of Science and Technology of Hanoi
}

\begin{document}

\maketitle

% ======================
\section{Introduction}
% ======================

Electrocardiogram (ECG) signals are one of the most commonly used biomedical
signals for monitoring cardiac activity and assessing the condition of the
cardiovascular system. By recording the electrical behavior of the heart,
ECG signals allow clinicians to identify abnormal heart rhythms and various
cardiac disorders such as arrhythmias. However, the continuous acquisition
of ECG data in both clinical and ambulatory settings produces a large volume
of signals, making manual analysis labor-intensive and prone to subjective
errors.

With the rapid development of machine learning and deep learning techniques,
automatic ECG signal analysis has gained increasing attention. These methods
are capable of learning discriminative patterns directly from data and can
assist medical professionals by providing reliable and consistent heartbeat
classification.

In this practical work, ECG heartbeat categorization is studied using a
publicly available dataset from Kaggle. Two different approaches are
implemented and compared: a traditional machine learning model based on
Random Forest, and a deep learning model using a one-dimensional
Convolutional Neural Network (1D-CNN). Special attention is paid to the
problem of class imbalance, which is a common issue in medical datasets.
The performance of the proposed methods is evaluated and compared with the
results reported in the original reference paper.

% ======================
\section{Dataset Overview}
% ======================

The dataset used in this study is the ECG Heartbeat Categorization Dataset
available on Kaggle. It is derived from the MIT-BIH Arrhythmia Database and
contains preprocessed ECG heartbeats extracted from raw recordings.

The dataset is divided into two subsets:
\begin{itemize}
    \item \textbf{Training set:} 87,554 samples.
    \item \textbf{Testing set:} 21,892 samples.
\end{itemize}

Each sample corresponds to a single ECG heartbeat represented as a
one-dimensional time-series vector of length 187. The last column of each
sample indicates the class label. The heartbeats are categorized into five
classes according to the AAMI standard:
\begin{enumerate}
    \item Normal
    \item Supraventricular
    \item Ventricular
    \item Fusion
    \item Unknown
\end{enumerate}

All ECG signals are normalized and padded to a fixed length, allowing them to
be directly used as inputs for machine learning and deep learning models.

% ======================
\section{Class Distribution}
% ======================

An important characteristic of the training dataset is the significant class
imbalance. The number of samples in each class is summarized as follows:
\begin{itemize}
    \item Normal: 72,471 samples
    \item Supraventricular: 2,223 samples
    \item Ventricular: 5,788 samples
    \item Fusion: 641 samples
    \item Unknown: 6,431 samples
\end{itemize}

As illustrated in Fig.~\ref{fig:class_dist}, the Normal class dominates the
dataset, while Supraventricular and Fusion beats are relatively rare. This
imbalance can negatively affect classification performance, especially for
minority classes, and therefore must be considered during model training and
evaluation.

\begin{figure}[htbp]
    \centering
    \includegraphics[width=0.65\linewidth]{figure/Distri.png}
    \caption{Class distribution of the training dataset, highlighting the
    imbalance between different heartbeat categories.}
    \label{fig:class_dist}
\end{figure}

% ======================
\section{Methodology}
% ======================

\subsection{Approach 1: Random Forest with GridSearchCV}

As a baseline machine learning method, a Random Forest classifier is
implemented. Random Forest is an ensemble-based algorithm that combines
multiple decision trees and is known for its robustness and good performance
on structured data.

To optimize the model, GridSearchCV is employed to search for the best
hyperparameters. In order to reduce the impact of class imbalance during
training, the weighted F1-score (\texttt{f1\_weighted}) is used as the
optimization metric. The optimal parameters obtained through this process
are:
\begin{itemize}
    \item Number of trees: 100
    \item Maximum depth: None
    \item Minimum samples split: 5
\end{itemize}

\subsection{Approach 2: 1D-CNN with SMOTE}

To further improve classification performance, especially for minority
classes, a deep learning approach based on a one-dimensional Convolutional
Neural Network (1D-CNN) is adopted. Since ECG signals are temporal sequences,
1D convolutions are well-suited for extracting local temporal features from
heartbeats.

Before training the CNN model, the Synthetic Minority Over-sampling
Technique (SMOTE) is applied to the training data in order to balance the
class distribution. By generating synthetic samples for underrepresented
classes, SMOTE helps the model learn more representative decision boundaries.

The CNN architecture consists of multiple convolutional and pooling layers,
followed by fully connected layers and a softmax output layer for
five-class classification.

% ======================
\section{Results}
% ======================

\subsection{Random Forest Performance}

The Random Forest model achieves a test accuracy of \textbf{97.42\%}.
Although the overall accuracy is high, the classification report reveals
that the model struggles to correctly identify minority classes. In
particular, the recall for Supraventricular and Fusion beats is relatively
low, which can be attributed to the strong class imbalance in the dataset.
\begin{table}[htbp]
\caption{Random Forest Classification Report}
\centering
\begin{tabular}{lcccc}
\hline
Class & Precision & Recall & F1-score & Support \\
\hline
Normal & 0.97 & 1.00 & 0.99 & 18118 \\
Supraventricular & 0.95 & 0.60 & 0.74 & 556 \\
Ventricular & 0.98 & 0.89 & 0.93 & 1448 \\
Fusion & 0.82 & 0.61 & 0.70 & 162 \\
Unknown & 1.00 & 0.95 & 0.97 & 1608 \\
\hline
Accuracy &  &  & 0.97 & 21892 \\
Weighted Avg & 0.97 & 0.97 & 0.97 & 21892 \\
\hline
\end{tabular}
\label{tab:rf_report}
\end{table}

The confusion matrix for the Random Forest model is shown in
Fig.~\ref{fig:rf_cm}, illustrating frequent misclassification of rare
heartbeat categories.

\begin{figure}[htbp]
    \centering
    \includegraphics[width=0.7\linewidth]{figure/rfs.png}
    \caption{Confusion matrix of the Random Forest model.}
    \label{fig:rf_cm}
\end{figure}

\subsection{CNN Performance with SMOTE}

After applying SMOTE and training the 1D-CNN model, the test accuracy
increases to \textbf{98.36\%}. More importantly, the recall for minority
classes improves significantly. For example, the recall of
Supraventricular beats increases from approximately 0.60 to 0.82, and the
recall of Fusion beats increases from approximately 0.61 to 0.83.
\begin{table}[htbp]
\caption{Random Forest Classification Report}
\centering
\begin{tabular}{lcccc}
\hline
Class & Precision & Recall & F1-score & Support \\
\hline
Normal & 0.97 & 1.00 & 0.99 & 18118 \\
Supraventricular & 0.95 & 0.60 & 0.74 & 556 \\
Ventricular & 0.98 & 0.89 & 0.93 & 1448 \\
Fusion & 0.82 & 0.61 & 0.70 & 162 \\
Unknown & 1.00 & 0.95 & 0.97 & 1608 \\
\hline
Accuracy &  &  & 0.97 & 21892 \\
Weighted Avg & 0.97 & 0.97 & 0.97 & 21892 \\
\hline
\end{tabular}
\label{tab:rf_report}
\end{table}

The confusion matrix of the CNN model, presented in
Fig.~\ref{fig:cnn_cm}, confirms that the CNN combined with SMOTE provides
better discrimination across all heartbeat categories.

\begin{figure}[htbp]
    \centering
    \includegraphics[width=0.7\linewidth]{figure/cnn.png}
    \caption{Confusion matrix of the CNN model trained with SMOTE.}
    \label{fig:cnn_cm}
\end{figure}

% ======================
\section{Conclusion}
% ======================

In this practical work, ECG heartbeat categorization is addressed using both
traditional machine learning and deep learning approaches. While the Random
Forest model achieves strong overall accuracy, it shows clear limitations in
detecting rare heartbeat classes due to data imbalance. By applying SMOTE
and using a 1D-CNN architecture, both the overall accuracy and the recall of
minority classes are significantly improved.

These results highlight the importance of handling class imbalance when
designing machine learning systems for medical applications and demonstrate
the effectiveness of deep learning methods for ECG signal classification.

\bibliographystyle{IEEEtran}
\end{document}
